% Compile with latexmk (uses ~/.latexmkrc: LuaLaTeX, output in build/)
% Week 8: combinatorial analysis, sample space, events.
% Sessions 1 and 2 are held on the SAME DAY, back to back (2 x 90 min,
% short break in between).
% Sources: lecture slides (Lectures 16--17); Ross, A First Course in
% Probability, 8th ed., Ch. 1--2 (problems paraphrased).
% The exercises are in exercises/ (one file each), placed by \exerciseframe.

\documentclass[aspectratio=169,11pt,usepdftitle=false]{beamer}
\usepackage{fontspec}
\setsansfont{Latin Modern Sans}
\usepackage{amsmath,amssymb,mathtools}
\usepackage{tikz}
\usetikzlibrary{calc}
\usepackage{booktabs}
\usepackage{cancel}

\usetheme{Madrid}

\definecolor{navy}{RGB}{24,48,75}
\definecolor{teal}{RGB}{0,120,125}
\definecolor{gold}{RGB}{190,130,20}

\setbeamercolor{structure}{fg=navy}
\setbeamercolor{alerted text}{fg=teal}
\setbeamercolor{block title}{bg=navy,fg=white}
\setbeamercolor{block body}{bg=navy!5}
\setbeamercolor{block title example}{bg=teal,fg=white}
\setbeamercolor{block body example}{bg=teal!6}
\setbeamercolor{block title alerted}{bg=gold,fg=white}
\setbeamercolor{block body alerted}{bg=gold!8}

\setbeamertemplate{navigation symbols}{}
\setbeamertemplate{footline}{%
  \usebeamerfont{footline}\hspace{1em}\insertshortauthor\hfill%
  \insertshorttitle\quad|\quad%
  \insertframenumber/\inserttotalframenumber%
  \hspace{1em}\vspace{2mm}%
}
\setbeamertemplate{itemize items}[circle]

% Contact (beamer loads hyperref itself)
\newcommand{\myemail}{tiborsandor2001@gmail.com}

\newcommand{\N}{\mathbb{N}}
\newcommand{\Om}{\Omega}
\newcommand{\sol}{\alert{Solution.}\ }

% Small Venn diagram of two events E, F in Omega; #1 = shading commands
\newcommand{\venn}[1]{%
  \begin{tikzpicture}[scale=.42,baseline=(O.base)]
    \node (O) at (.5,0) {\phantom{x}};
    #1
    \draw[thick] (-1.5,-1.15) rectangle (2.5,1.15);
    \draw[thick,navy] (0,0) circle (.85);
    \draw[thick,navy] (1,0) circle (.85);
    \node at (-.35,0) {\scriptsize $E$};
    \node at (1.35,0) {\scriptsize $F$};
    \node at (2.15,.8) {\scriptsize $\Omega$};
  \end{tikzpicture}}
\newcommand{\shadeall}[1]{\fill[#1] (-1.5,-1.15) rectangle (2.5,1.15);}

% The exercises (one file each in exercises/) and \exerciseframe
% Exercises of Week 8, shared by beamer.tex, the study guide, the script and
% exercises.tex. Every exercise lives in a file of its own,
% exercises/NN-slug.tex, numbered 8.NN (see exercises.lua for the docstring).
% The body of an exercise file:
%
%   \begin{problem} ... \end{problem}      the statement, used everywhere;
%                                          inside it, \begin{context} ...
%                                          \end{context} is background that
%                                          the slides leave out
%   \begin{slideproblem} ... \end{slideproblem}
%                                          (optional) a shorter statement
%                                          for the slides
%   \begin{slides} ... \end{slides}        one frame of the deck: \exstatement
%                                          places the statement; repeat the
%                                          environment for more frames
%   \begin{solution} ... \end{solution}    the detailed solution (study
%                                          guide, exercise PDFs, website)
%
% Commands for the documents (the argument is the slug: 01-license-plates.tex
% has the slug license-plates):
%
%   \exref{slug}                 the number, e.g. 8.1, for "Exercise~\exref{...}"
%   \extitle{slug}               the title from the docstring
%   \exerciseframe[opts]{slug}   beamer: the frames of the exercise
%   \studyexercise{slug}         study guide: heading, statement, solution
%   \scriptexercise{slug}        script: the statement box, as on the slide
%   \allexercises                exercises.tex: every exercise, in order; on
%                                the website statement and solution with
%                                markers where the pipeline splits the HTML
%                                (\collectionexercise), in the PDF plain:
%                                each statement, followed by its solution if
%                                \ifsolutions (exercises-solutions.tex)

\def\exweek{8}

\directlua{exercises = dofile("exercises/exercises.lua")
  exercises.scan("exercises", "\exweek")}

\ExplSyntaxOn

\msg_new:nnn { exercises } { unknown } { Unknown~exercise~`#1'. }

\cs_new_protected:Npn \__ex_missing:n #1
  { \msg_error:nnn { exercises } { unknown } {#1} ?? }
\cs_new:Npn \exref #1
  {
    \cs_if_exist:cTF { exnum@#1 }
      { \use:c { exnum@#1 } } { \__ex_missing:n {#1} }
  }
\cs_new:Npn \extitle #1
  {
    \cs_if_exist:cTF { extitle@#1 }
      { \use:c { extitle@#1 } } { \__ex_missing:n {#1} }
  }

% ---------------------------------------------------------------------------
% Loading an exercise: its parts are stored, nothing is typeset. The slides
% typeset the stored parts; the other documents read the part they need from
% the file again (\__ex_input_part:nn), as stored tables do not survive lwarp.

\tl_new:N \g__ex_current_tl
\tl_new:N \g__ex_problem_tl
\tl_new:N \g__ex_slideproblem_tl
\tl_new:N \g__ex_solution_tl
\seq_new:N \g__ex_slides_seq
\box_new:N \l__ex_scratch_box

\cs_new_protected:Npn \__ex_load:n #1
  {
    \cs_if_exist:cF { exfile@#1 } { \__ex_missing:n {#1} }
    \tl_gset:Nn \g__ex_current_tl {#1}
    \tl_gclear:N \g__ex_problem_tl
    \tl_gclear:N \g__ex_slideproblem_tl
    \tl_gclear:N \g__ex_solution_tl
    \seq_gclear:N \g__ex_slides_seq
    \vbox_set:Nn \l__ex_scratch_box
      {
        \DeclareDocumentEnvironment { problem } { +b }
          { \tl_gset:Nn \g__ex_problem_tl {##1} } { }
        \DeclareDocumentEnvironment { slideproblem } { +b }
          { \tl_gset:Nn \g__ex_slideproblem_tl {##1} } { }
        \DeclareDocumentEnvironment { slides } { +b }
          { \seq_gput_right:Nn \g__ex_slides_seq {##1} } { }
        \DeclareDocumentEnvironment { solution } { +b }
          { \tl_gset:Nn \g__ex_solution_tl {##1} } { }
        \exp_args:Nv \input { exfile@#1 }
      }
  }

% Typesets one part (problem, slideproblem, solution) of an exercise straight
% from its file; the other parts are skipped.
\cs_new_protected:Npn \__ex_input_part:nn #1#2
  {
    \group_begin:
      \clist_map_inline:nn { problem, slideproblem, slides, solution }
        {
          \str_if_eq:nnTF {##1} {#2}
            { \DeclareDocumentEnvironment {##1} { } { } { } }
            { \DeclareDocumentEnvironment {##1} { +b } { } { } }
        }
      \exp_args:Nv \input { exfile@#1 }
    \group_end:
  }

% Exercise 8.1, or Exercise 8.2 · Finance
\cs_new:Npn \__ex_heading:
  {
    Exercise~\exref { \g__ex_current_tl }
    \cs_if_exist:cT { exfinance@ \g__ex_current_tl } { ~·~Finance }
  }
\prg_new_conditional:Npnn \__ex_if_finance: { T, F, TF }
  {
    \cs_if_exist:cTF { exfinance@ \g__ex_current_tl }
      { \prg_return_true: } { \prg_return_false: }
  }

% ---------------------------------------------------------------------------
% Beamer

% The finance context is printed in the study guide, the exercise PDFs and on
% the website; the
% slides and the script (which shows what the slide shows) leave it out.
\bool_new:N \l__ex_context_bool
\bool_set_true:N \l__ex_context_bool
\DeclareDocumentEnvironment { context } { }
  {
    \bool_if:NTF \l__ex_context_bool
      { \par \textit { Finance~context. } ~ }
      { \__ex_gobble_context:w }
  }
  { \bool_if:NT \l__ex_context_bool { \par } }
\cs_new_protected:Npn \__ex_gobble_context:w #1 \end #2
  { \str_if_eq:nnTF {#2} { context } { \end { context } } { \__ex_gobble_context:w } }

\IfClassLoadedTF { beamer }
  {
    \bool_set_false:N \l__ex_context_bool

    \cs_new:Npn \__ex_block: { \__ex_if_finance:TF { exampleblock } { block } }
    \NewDocumentCommand \exstatement { }
      {
        \begin { \__ex_block: } { \__ex_heading: }
          \tl_if_empty:NTF \g__ex_slideproblem_tl
            { \g__ex_problem_tl } { \g__ex_slideproblem_tl }
        \end { \__ex_block: }
      }

    \tl_new:N \l__ex_frames_tl
    \NewDocumentCommand \exerciseframe { O{t} m }
      {
        \__ex_load:n {#2}
        \tl_clear:N \l__ex_frames_tl
        \seq_map_indexed_inline:Nn \g__ex_slides_seq
          {
            \tl_put_right:Nn \l__ex_frames_tl { \begin{frame}[#1] }
            \tl_put_right:Ne \l__ex_frames_tl
              {
                {
                  \exp_not:N \extitle {#2}
                  \int_compare:nNnT {##1} > { 1 } { \c_space_tl (cont.) }
                }
              }
            \tl_put_right:Nn \l__ex_frames_tl { ##2 \end{frame} }
          }
        \tl_use:N \l__ex_frames_tl
      }
  }
  {
    \cs_new:Npn \__ex_box: { \__ex_if_finance:TF { fin } { task } }
    \cs_new_protected:Npn \__ex_statement_box:n #1
      {
        \begin { \__ex_box: } { \__ex_heading: }
          \__ex_input_part:nn { \g__ex_current_tl } {#1}
        \end { \__ex_box: }
      }

    \NewDocumentCommand \studyexercise { m }
      {
        \__ex_load:n {#1}
        \subsection* { Exercise~\exref{#1}~---~\extitle{#1} }
        \addcontentsline { toc } { subsection }
          { Exercise~\exref{#1}~---~\extitle{#1} }
        \__ex_statement_box:n { problem }
        \soltitle \__ex_input_part:nn {#1} { solution }
      }

    \NewDocumentCommand \scriptexercise { m }
      {
        \__ex_load:n {#1}
        \group_begin:
          \bool_set_false:N \l__ex_context_bool
          \tl_if_empty:NTF \g__ex_slideproblem_tl
            { \__ex_statement_box:n { problem } }
            { \__ex_statement_box:n { slideproblem } }
        \group_end:
      }

    % The website (lwarp): statement and solution of each exercise, between
    % markers where the pipeline splits the HTML.
    \NewDocumentCommand \collectionexercise { m }
      {
        \__ex_load:n {#1}
        \par [[exercise:#1:problem]] \par
        \__ex_input_part:nn {#1} { problem }
        \par [[exercise:#1:solution]] \par
        \__ex_input_part:nn {#1} { solution }
        \par [[exercise:#1:end]] \par
      }
    % The PDF: plain, each statement, followed by its solution in the
    % version with solutions (\ifsolutions). The statement is set off like a
    % quote in Markdown, by a bar on the left (quote=<color>, also used by
    % exercises.tex) in the color of its box in the study guide: navy, or teal
    % for a finance exercise.
    \tcbset
      {
        quote/.style =
          {
            enhanced, breakable, frame~hidden, interior~hidden,
            sharp~corners, boxrule = 0pt, borderline~west = {2pt}{0pt}{#1},
            left = 8pt, right = 0pt, top = 1pt, bottom = 1pt,
            before~skip = 6pt, after~skip = 6pt,
            % the document's paragraph spacing (parskip=half)
            before~upper = { \setlength { \parskip } { .5\baselineskip } }
          }
      }
    \NewDocumentCommand \plainexercise { m }
      {
        \__ex_load:n {#1}
        \subsection* { Exercise~\exref{#1}~---~\extitle{#1} }
        \addcontentsline { toc } { subsection }
          { Exercise~\exref{#1}~---~\extitle{#1} }
        \__ex_if_finance:TF
          { \begin{tcolorbox}[quote=teal] }
          { \begin{tcolorbox}[quote=navy] }
          \__ex_input_part:nn {#1} { problem }
        \end{tcolorbox}
        \legacy_if:nT { solutions }
          {
            \soltitle
            \__ex_input_part:nn {#1} { solution }
          }
      }
    \NewDocumentCommand \allexercises { }
      {
        \IfPackageLoadedTF { lwarp }
          { \directlua { exercises.each("collectionexercise") } }
          { \directlua { exercises.each("plainexercise") } }
      }
  }

\ExplSyntaxOff


% PDF properties and the date on the title slide
% The properties of every PDF of Week 8: author (with email address), title,
% subject, keywords and the date the document last changed, in the PDF
% (hyperref) and on its title page (\date), and the header and footer of the
% printed documents (below). Each document loads this file after hyperref and
% then calls \documentinfo{<what it is>}, e.g. \documentinfo{Slides}.
%
% The date: the website build (scripts/build-content.mjs) puts
% \def\lastmodified{YYYY-MM-DD} in front of the document, the date of the last
% commit that changed its sources. Compiled locally, it is today.

\newcommand{\courseinfo}{Mathematics for Finance}
\newcommand{\weeklabel}{Week 8}
\newcommand{\weekinfo}{\weeklabel: Combinatorial analysis, sample space, events}
\newcommand{\authorname}{Tibor SÁNDOR}
\providecommand{\myemail}{tiborsandor2001@gmail.com}

\makeatletter
\ifdefined\lastmodified\else
	\edef\lastmodified{\the\year-\two@digits{\the\month}-\two@digits{\the\day}}
\fi
\def\di@split#1-#2-#3\@nil{%
	\def\di@year{#1}\def\di@month{#2}\def\di@day{#3}}
\expandafter\di@split\lastmodified\@nil
% November 11, 2026
\edef\lastmodifiedtext{%
	\ifcase\numexpr\di@month\relax\or January\or February\or March\or April\or
		May\or June\or July\or August\or September\or October\or November\or
		December\fi\space\number\di@day, \di@year}
% D:20261111000000, the date format of PDF metadata
\edef\di@pdfdate{D:\di@year\di@month\di@day000000}

\newcommand{\documentinfo}[1]{%
	\hypersetup{
		pdftitle={\courseinfo, \weekinfo\space -- #1},
		pdfauthor={\authorname\space<\myemail>},
		pdfsubject={\courseinfo, \weekinfo},
		pdfkeywords={probability, combinatorics, mathematics for finance,
				Corvinus University of Budapest},
		pdfcreationdate={\di@pdfdate},
		pdfmoddate={\di@pdfdate},
		pdflang={en-US},
	}%
	\date{Last updated: \lastmodifiedtext}%
	\gdef\documentkind{#1}%
}

% Header and footer of the printed documents (KOMA-Script, scrlayer-scrpage;
% the slides have beamer's footline, the website neither): the course, the
% week and the document on the left of the header (the week's topic is on the
% title page), the university on the right; in the footer the author and the
% page as <current>/<total>. The title page gets only the footer.
\newcommand{\documentkind}{}
\IfClassLoadedTF{beamer}{}{\IfPackageLoadedTF{lwarp}{}{%
	\usepackage{lastpage}%
	\usepackage[headsepline=1.5pt,footsepline=1.5pt,plainfootsepline]
		{scrlayer-scrpage}%
	\clearpairofpagestyles
	\setkomafont{pageheadfoot}{\small\sffamily\color{navy}}%
	\setkomafont{pagenumber}{}%
	\setkomafont{headsepline}{\color{navy!35}}%
	\setkomafont{footsepline}{\color{navy!35}}%
	\ihead{\courseinfo\ \textperiodcentered\ \weeklabel\ \textperiodcentered\ 
		\documentkind}%
	\ohead{Corvinus University of Budapest}%
	\ifoot*{\authorname\ \textperiodcentered\ \myemail}%
	\ofoot*{\pagemark/\pageref*{LastPage}}%
	\pagestyle{scrheadings}%
	% Extra stretch for breaking lines with long inline formulas.
	\setlength{\emergencystretch}{1.5em}%
}}
\makeatother

\documentinfo{Slides}

\title[Mathematics for Finance]{Mathematics for Finance}
\author{Tibor SÁNDOR}
\subtitle{Week 8: Combinatorial analysis, sample space, events}
\institute{Corvinus University of Budapest\\[2pt]
  \href{mailto:\myemail}{\texttt{\myemail}}}

\begin{document}

\begin{frame}
	\titlepage
\end{frame}

\begin{frame}{Plan for today}
	\begin{enumerate}
		\item \textbf{Session 1 (90 min):} counting principle; permutations;
		      permutations with repetition; multinomial coefficients;
		      combinations; lattice paths.
		\item \textbf{Session 2 (90 min):} Pascal's triangle; the binomial tree;
		      binomial and multinomial theorems; stars and bars;
		      sample space and events; De Morgan's laws; inclusion--exclusion.
	\end{enumerate}
	\vspace{.3cm}
	\textbf{Contact:} Tibor SÁNDOR,
	\href{mailto:\myemail}{\texttt{\myemail}}
\end{frame}

%======================================================================
\section{Session 1: Counting}
%======================================================================

\begin{frame}
	\centering\Large\color{navy}Session 1: Permutations and combinations
\end{frame}

\begin{frame}{Toolbox}
	\begin{block}{Basic principle of counting}
		\centering
		$m$ ways, then $n$ ways for each of them
		$\;\Longrightarrow\;$ \alert{$m\cdot n$} ways together
	\end{block}
	% \begin{columns}[T]
	% 	\begin{column}{.48\textwidth}
	% 		\begin{block}{Permutations \hfill\normalfont\small order matters}
	% 			\centering
	% 			$n$ objects in a row: \ \alert{$n!$}\\[3mm]
	% 			$k$ of $n$, in order:\\[1mm]
	% 			\alert{$\dfrac{n!}{(n-k)!}$} $=n(n-1)\cdots(n-k+1)$
	% 		\end{block}
	% 	\end{column}
	% 	\begin{column}{.48\textwidth}
	% 		\begin{block}{Combinations \hfill\normalfont\small order does not matter}
	% 			\centering
	% 			$k$ of $n$, as a set:\\[1mm]
	% 			\alert{$\dbinom nk$} $=\dfrac{n!}{k!\,(n-k)!}$\\[2mm]
	% 			{\small each set comes in $k!$ orders}
	% 		\end{block}
	% 	\end{column}
	% \end{columns}
	\begin{block}{Two questions to ask every time}
		\centering
		\textbf{1.} Does the order matter? \qquad
		\textbf{2.} Can an element be used more than once?
		\begin{center}
			\begin{tabular}{lcc}\toprule
				              & ordered              & unordered                       \\\midrule
				no repetition & $\dfrac{n!}{(n-k)!}$ & $\displaystyle\binom nk$        \\[5mm]
				repetition    & $n^k$                & $\displaystyle\binom{n+k-1}{k}$ \\\bottomrule
			\end{tabular}
		\end{center}
	\end{block}
\end{frame}

\exerciseframe{license-plates}
\exerciseframe{couples-at-a-concert}
\exerciseframe{bus-queue}
\exerciseframe{letter-arrangements}
\exerciseframe{eight-digit-numbers}

\begin{frame}{Multinomial coefficient}
	\begin{block}{Multinomial coefficient}
		Number of ways to divide $n$ distinct objects into $r$ labeled groups of
		sizes $n_1,\dots,n_r$ ($n_1+\dots+n_r=n$):
		\[ \binom{n}{n_1,\dots,n_r} = \frac{n!}{n_1!\cdots n_r!}
			= \binom{n}{n_1}\binom{n-n_1}{n_2}\cdots\binom{n_r}{n_r}. \]
	\end{block}
	The same formula as for identical objects in a row: a division into groups
	is a word with $n_1$ letters ``1'', $n_2$ letters ``2'', \dots
\end{frame}

\exerciseframe{dividing-into-groups}
\exerciseframe{rating-bonds}
\exerciseframe{up-and-down-days}
\exerciseframe{committees}
\exerciseframe{lattice-paths}

\begin{frame}{Session 1: summary}
	\small
	\begin{itemize}
		\item Ordered, no repetition: $n!/(n-k)!$; ordered, repetition: $n^k$.
		\item Unordered, no repetition: $\binom nk$ (e.g.\ lattice paths).
		\item Identical objects in a row, or groups of given sizes: the
		      multinomial coefficient $n!/(n_1!\cdots n_r!)$.
		\item ``Together'' $\to$ block; ``not together'' $\to$ complement.
	\end{itemize}
\end{frame}

%======================================================================
\section{Session 2: Pascal's triangle, events}
%======================================================================

\begin{frame}
	\centering\Large\color{navy}Session 2:
	Pascal's triangle, sample space, events
\end{frame}

\begin{frame}{Recap of Session 1}
	\begin{itemize}
		\item Ask: \textbf{does order matter?} \ \textbf{repetition allowed?}
		      \begin{center}
			      \begin{tabular}{lcc}\toprule
				                    & ordered              & unordered                       \\\midrule
				      no repetition & $\dfrac{n!}{(n-k)!}$ & $\displaystyle\binom nk$        \\[5mm]
				      repetition    & $n^k$                & $\displaystyle\binom{n+k-1}{k}$ \\\bottomrule
			      \end{tabular}
		      \end{center}
		\item Identical objects in a row, or groups of given sizes:
		      $n!/(n_1!\cdots n_r!)$.
		\item ``Together'' $\to$ block; ``not together'' or ``at least one''
		      $\to$ complement.
		\item Lattice paths: choose which of the $n$ moves go up, $\binom nk$.
	\end{itemize}
	\vspace{2mm}
	\textbf{Now:} Pascal's triangle, the binomial theorem, identical objects
	into boxes, and the language of events.
\end{frame}

\begin{frame}[t]{Pascal's triangle}
	\begin{columns}[T]
		\begin{column}{.52\textwidth}
			\centering
			\begin{tikzpicture}[every node/.style={font=\small}]
				\foreach \n in {0,...,6} {
						\node[navy!70] at (-3.9,-.62*\n) {\scriptsize $n=\n$};
						\foreach \k in {0,...,\n} {
								\pgfmathtruncatemacro{\c}{factorial(\n)/(factorial(\k)*factorial(\n-\k))}
								\node[inner sep=1pt,minimum size=13pt] (p\n-\k)
								at ({(\k-\n/2)*.92},{-.62*\n}) {\c};
							}
					}
				\draw[gold,thick] (p4-1) circle (7pt);
				\draw[gold,thick] (p4-2) circle (7pt);
				\draw[teal,thick,fill=teal!15,fill opacity=.5] (p5-2) circle (7pt);
				\draw[gold,thick,->] (p4-1)--(p5-2);
				\draw[gold,thick,->] (p4-2)--(p5-2);
				\node[teal] at (p5-2) {\small 10};
			\end{tikzpicture}\\[1mm]
			\footnotesize row $n$, position $k=0,1,\dots,n$ holds $\binom nk$\\[3mm]
			\small $(a+b)^5=a^5+5a^4b+10a^3b^2+10a^2b^3+5ab^4+b^5$
		\end{column}
		\begin{column}{.46\textwidth}
			\small
			\begin{block}{Pascal's rule (proof: Exercise~\exref{pascals-rule})}
				$\displaystyle\binom nk+\binom n{k+1}=\binom{n+1}{k+1}$\\[1mm]
				builds the triangle: each entry is the sum of the two above,
				$\binom41+\binom42=4+6=10=\binom52$.
			\end{block}
			\textbf{Readings of row $n$:}
			\begin{itemize}
				\item paths with $k$ up-moves among $n$: lattice paths
				      (Exercise~\exref{lattice-paths}), the binomial tree
				      (Exercise~\exref{binomial-tree});
				\item coefficients of $(a+b)^n$ (see row 5);
				\item row sum $=2^n$; symmetry $\binom nk=\binom n{n-k}$.
			\end{itemize}
		\end{column}
	\end{columns}
\end{frame}

\exerciseframe{pascals-rule}

\exerciseframe{binomial-tree}

\begin{frame}{Binomial and multinomial theorems}
	\begin{block}{Expanding a power of a sum}
		\[ (a+b)^n=\sum_{k=0}^n\binom nk a^kb^{n-k},\qquad
			(x_1+\dots+x_r)^n=\sum_{n_1+\dots+n_r=n}\binom{n}{n_1,\dots,n_r}
			x_1^{n_1}\cdots x_r^{n_r}. \]
	\end{block}
	The coefficients of $(a+b)^n$ are row $n$ of Pascal's triangle; the
	multinomial coefficients count the ways to pick $x_1$ from $n_1$ of the
	$n$ factors, $x_2$ from $n_2$ of them, \dots
\end{frame}

\exerciseframe{expansions}

\begin{frame}{Stars and bars}
	\small
	\textbf{Question:} in how many ways can $n$ identical objects be put into
	$r$ distinct boxes?

	\textbf{Example:} 5 identical coins for 3 children. Write a star for each
	coin and a bar between two children:
	\begin{center}
		\begin{tabular}{c@{\quad means\quad}c@{\hspace{12mm}}c@{\quad means\quad}c}
			\texttt{**|*|**} & $(2,1,2)$ & \texttt{*****||} & $(5,0,0)$ \\
			\texttt{|***|**} & $(0,3,2)$ & \texttt{*||****} & $(1,0,4)$
		\end{tabular}
	\end{center}
	\vspace{-2mm}\pause
	\begin{itemize}
		\item Every division is a row of $5$ stars and $2$ bars, and every such
		      row is a division.
		\item The row has $7$ places; it is fixed by the places of the 2 bars:
		      \alert{$\binom72=21$} divisions.
	\end{itemize}
	\pause
	\begin{block}{Combinations with repetition (stars and bars)}
		$n$ identical objects into $r$ distinct boxes: $n$ stars and $r-1$ bars,
		\alert{$\binom{n+r-1}{r-1}$} ways.\\
		Every box nonempty: first put one object into each box, then divide the
		other $n-r$ freely: $\binom{n-1}{r-1}$ ways.
	\end{block}
	This is the ``unordered, with repetition'' entry of the table: choosing
	$k$ times from $n$ kinds is $k$ stars and $n-1$ bars, $\binom{n+k-1}{k}$.
\end{frame}

\exerciseframe{stars-and-bars}

\exerciseframe{identical-dice}

\exerciseframe{allocating-capital}

%----------------------------------------------------------------------
\begin{frame}{Sample space and events}
	\begin{block}{Definitions}
		\begin{itemize}
			\item \textbf{Sample space} $\Om$: set of all possible outcomes.
			\item \textbf{Event}: a subset $E\subseteq\Om$; $E$ \emph{occurs} if the
			      outcome lies in $E$.
			\item $E\cup F$ (or), \ $E\cap F=EF$ (and), \ $E^c$ (not),
			      \ $E\subseteq F$ ($E$ implies $F$).
			\item $E,F$ \textbf{mutually exclusive} if $EF=\emptyset$.
		\end{itemize}
	\end{block}
	\begin{center}
		\renewcommand{\arraystretch}{1}
		\setlength{\tabcolsep}{6mm}
		\begin{tabular}{ccc}
			\venn{\fill[teal!35] (0,0) circle (.85);}
			          & \venn{\fill[teal!35] (1,0) circle (.85);}
			          & \venn{\fill[teal!35] (0,0) circle (.85); \fill[white] (1,0) circle (.85);}                       \\
			$E$       & $F$                                                                        & $E\setminus F=EF^c$ \\[2mm]
			\venn{\fill[teal!35] (0,0) circle (.85); \fill[teal!35] (1,0) circle (.85);}
			          & \venn{\begin{scope}\clip (0,0) circle (.85);
					            \fill[teal!35] (1,0) circle (.85);\end{scope}}
			          & \venn{\shadeall{teal!35}\fill[white] (0,0) circle (.85);}                                        \\
			$E\cup F$ & $EF$                                                                       & $E^c$
		\end{tabular}
	\end{center}
\end{frame}

\begin{frame}[t]{De Morgan's laws for two events}
	\begin{block}{De Morgan's laws}
		\[ (E\cup F)^c = E^c\cap F^c, \qquad (E\cap F)^c = E^c\cup F^c. \]
		``not (E or F)'' $=$ ``neither E nor F'';\quad
		``not (both E and F)'' $=$ ``at least one of them fails''.
	\end{block}
	\vspace{1mm}
	\begin{center}
		\renewcommand{\arraystretch}{1.3}
		\begin{tabular}{ccccc}
			\venn{\shadeall{teal!35}\fill[white] (0,0) circle (.85);}
			                  & $\cap$ &
			\venn{\shadeall{gold!40}\fill[white] (1,0) circle (.85);}
			                  & $=$    &
			\venn{\shadeall{navy!35}\fill[white] (0,0) circle (.85);
				\fill[white] (1,0) circle (.85);}       \\[-1mm]
			\scriptsize $E^c$ &        & \scriptsize $F^c$ &  &
			\scriptsize $E^c\cap F^c$: outside both $=(E\cup F)^c$        \\
			\venn{\shadeall{teal!35}\fill[white] (0,0) circle (.85);}
			                  & $\cup$ &
			\venn{\shadeall{gold!40}\fill[white] (1,0) circle (.85);}
			                  & $=$    &
			\venn{\shadeall{navy!35}\begin{scope}\clip (0,0) circle (.85);
					\fill[white] (1,0) circle (.85);\end{scope}} \\[-1mm]
			\scriptsize $E^c$ &        & \scriptsize $F^c$ &  &
			\scriptsize $E^c\cup F^c$: all but the lens $=(E\cap F)^c$
		\end{tabular}
	\end{center}
	\vspace{-1mm}\small Shade $E^c$ and $F^c$: shaded \emph{twice} $=$
	intersection; shaded \emph{at least once} $=$ union.
\end{frame}

\begin{frame}[t]{De Morgan's laws: general form}
	\begin{block}{General form ($n$ events)}
		\[ \Big(\bigcup_{i=1}^n E_i\Big)^c = \bigcap_{i=1}^n E_i^c, \qquad
			\Big(\bigcap_{i=1}^n E_i\Big)^c = \bigcup_{i=1}^n E_i^c. \]
	\end{block}
	Next week we use exactly this: $P(A^cB^c)=1-P(A\cup B)$.
\end{frame}

\exerciseframe{sample-space}

\exerciseframe{stock-path-events}

\exerciseframe{credit-basket}

\exerciseframe{bank-products}

\begin{frame}{Session 2: summary}
	\small
	\begin{itemize}
		\item Binomial tree: $\binom nk$ of the $2^n$ paths have $k$ up-moves;
		      Pascal's rule arranges these numbers in Pascal's triangle.
		\item Binomial and multinomial theorems: the same coefficients appear
		      in $(a+b)^n$ and $(x_1+\dots+x_r)^n$.
		\item Identical objects into boxes: stars and bars.
		\item Events are sets: $\cup$ = or, $\cap$ = and, $^c$ = not;
		      De Morgan: $(E\cup F)^c=E^cF^c$, $(EF)^c=E^c\cup F^c$;
		      inclusion--exclusion with a Venn diagram, from the inside out.
	\end{itemize}
\end{frame}

\end{document}
